%% kapa: exam=1 round=22 year=2011 period=1 form=A subject=civil_law subject_ko=민법 seq=22 number=22
%% source: Q-Net 감정평가사 기출문제 (cst00309) artlSeq=5036790 file=extracted/1008/5036790/2011년 제22회 감정평가사 1차시험 기출문제/1교시.hwp
%% conversion: kordoc 4.24.51 -> tools/gen1.py (auto); review=proofread-sonnet; flags=none
\begin{question}{22}
甲은 乙에게 인감도장을 주고 자신의 아파트에 관한 임대 등 일체의 관리를 위임하였고, 丙은 甲으로 행세하는 乙로부터 아파트를 임차하여 임차기간을 갱신하면서 거주하여 오다가 그 아파트를 매수하였다. 乙은 임대차계약 당시뿐만 아니라 매매계약 당시에도 甲의 인감도장을 사용하여 甲의 이름으로 계약을 체결하고, 그 대금 등도 甲의 이름으로 수령하였다. 다음 설명으로 옳은 것은? (다툼이 있으면 판례에 의함)
\begin{choices}
  \item 乙의 행위는 진의 아닌 의사표시에 해당하므로, 丙이 乙의 진의 아님을 알았거나 알 수 있었다면 무효로 된다.
  \item 乙의 행위는 통정허위표시에 해당하므로, 甲은 매매계약의 무효를 주장할 수 있다.
  \item 乙이 甲으로 가장하여 아파트를 임대하였으므로, 임대차계약은 甲에게 효력이 발생하지 않는다.
  \item 丙은 甲에게 아파트 소유권이전등기를 청구할 수 있다.
  \item 乙은 甲에게 매매계약을 추인할 것인지 여부의 확답을 최고할 수 있다.
\end{choices}
\end{question}
\answer{4}
